% Interim short version of the Method chapter for the supervisor draft of
% 29 July 2026 (compiled via BA_Zwischenstand_2026-07-29.tex, NOT via main.tex).
% The full chapter will be written in sections/method.tex; every number here is
% frozen method documentation (see CLAUDE.md), not derived from the coding table.
\chapter{Method}
\label{sec:method-summary}

\noindent\emph{Interim summary. The full chapter --- search documentation, screening protocol, coding procedure, and PRISMA flow diagram --- follows in the next version of this draft.}

\bigskip

\noindent The review follows the protocol announced in the proposal. The documented Scopus query was run on 17 July 2026 and returned 432 unique records: peer-reviewed journal articles in English, published 2015--2026, in the Scopus subject areas business, economics, and decision sciences. The query finds all three benchmark studies defined in advance (Babina et al., 2024; Krakowski et al., 2023; Lui et al., 2022).

For journal quality, the strict AJG~4*/4 threshold of the proposal was replaced by a broader, uniform inclusion rule: a record stays in when its journal is listed in the AJG~2024 with a rating of at least~2. One fixed ranking edition applies to all records, so the rule is transparent and non-selective. This filter reduced the corpus from 432 to 174 records; most of the excluded records appeared in journals not listed in the ranking.

Title and abstract screening tested four criteria: the study is empirical, works at the firm level, treats firms' AI investment or adoption as its object, and measures a performance or competitive-advantage outcome. This step reduced the 174 records to 81. Reading the full texts removed another fourteen studies, twelve on content grounds and two that could not be retrieved. The final corpus holds 67 studies. Every screening decision is recorded with a reason code, and the full chain (432 $\rightarrow$ 174 $\rightarrow$ 81 $\rightarrow$ 67) will be reported as a PRISMA-style flow diagram.

Each included study was coded into the state-of-the-art table announced in the proposal; an excerpt is attached in the Appendix (Tables~\ref{tab:coding-a1} and~\ref{tab:coding-a2}). Following the supervisor's feedback, the scheme separates superior firm performance from competitive advantage as distinct outcome constructs, each with its own measurement column.
